\chapter{Considerações Finais}
\label{chp:conclusion}

Este trabalho tem como pergunta central antecipar, com algum rigor estatístico, se uma troca realizada por uma equipe da NBA durante a temporada resultaria em ganho de desempenho. A resposta para essa pergunta, construída ao longo dos quatro capítulos anteriores, não se reduz a uma resposta binária, tampouco poderia fazê-lo. O modelo multidimensional proposto na Seção~\ref{sec:objectives} decompôs deliberadamente esse problema em quatro hipóteses menores, cada uma testando uma dimensão diferente da troca, e cada uma recebeu sua própria resposta, relatada em detalhe no Capítulo 4 e resumida na Tabela~\ref{tab:sintese_resultados}. O que essas quatro respostas têm em comum, lidas em conjunto, é a magnitude, pois mesmo onde há suporte estatístico, o efeito é sempre fraco pela convenção de \citet{cohen1988}. As quatro variáveis testadas foram a desigualdade salarial, o ajuste funcional, a experiência dos jogadores adquiridos e a estrutura da negociação; nenhuma delas decide, isoladamente, se uma troca vai dar certo. Isso não invalida o modelo; é, antes, a confirmação empírica de uma premissa que já estava na Introdução, a de que a troca é um evento multidimensional demais para caber em um único indicador.

A dimensão salarial foi a que recebeu o tratamento mais extenso neste trabalho, e por um bom motivo: era a única das quatro em que a literatura prévia estava genuinamente dividida entre duas teorias com previsões opostas. O resultado encontrado foi uma associação negativa e estatisticamente significativa entre a concentração da folha salarial em quadra e o desempenho subsequente da equipe, o que favorece a Teoria da Equidade de \citet{adams1965} sobre a Teoria do Torneio de \citet{lazear1981}, na linha do que \citet{frick2003} relatou favorável ao torneio em outro contexto, mas na direção oposta. É um resultado real, robusto a diferentes larguras de janela de troca (Seção~\ref{sec:r_h1_main}) e confirmado por uma métrica de desigualdade escolhida antes de qualquer teste de correlação, para evitar que a métrica fosse ajustada ao resultado que ela própria viria a produzir (Seção~\ref{sec:r_h1_metric}). Mas é também um resultado que explica pouco: cerca de 3,6\% da variação de desempenho, e concentrado, como mostra a Subseção~\ref{sec:r_h1_moderation}, especificamente nas equipes de perfil competitivo Ofensivo. Esse recorte ainda não havia sido testado pela literatura revisada no Capítulo 2, já que \citet{katayama2011} mediu heterogeneidade pela probabilidade de acesso aos \textit{playoffs}, não pelo estilo de jogo do elenco. Equidade salarial parece importar, mas não para todo tipo de time da mesma forma, e este trabalho só consegue apontar essa concentração, não explicá-la.

A dimensão que menos se confirmou foi justamente aquela em que a metodologia da Contribuição Relativa de \citet{sa2025} havia sido adaptada com mais cuidado, a saber, o ajuste funcional do jogador adquirido. Times que trocam para preencher uma lacuna estrutural diagnosticável não superam, de forma estatisticamente distinguível de acaso, os times que reforçam uma posição já bem servida. Um sinal mais fraco da mesma relação sobrevive quando a magnitude da lacuna é tratada de forma contínua em vez de um corte binário, mas mesmo esse sinal contínuo só se mostra estável na janela 2019-2021, exatamente a única em que a fonte de dado posicional é a mesma usada por \citet{sa2025} originalmente (Seção~\ref{sec:r_h2}). Este trabalho não chegou a uma estratégia analítica capaz de isolar em que medida essa instabilidade decorre de uma limitação real da métrica ou da qualidade desigual do dado posicional entre eras — uma questão que permanece em aberto e é retomada na Seção~\ref{sec:future_work}.

As duas hipóteses restantes, referentes à experiência dos jogadores adquiridos e à estrutura do pacote de troca, encontraram algum suporte nos dados, ainda que qualificado por particularidades relevantes de cada resultado. Trocas por jogadores mais experientes tendem a se associar a um impacto de curto prazo mais positivo, um achado consistente com a literatura sobre curva de adaptação revisada no Capítulo 2, embora o efeito só se torne realmente perceptível no quartil mais velho da amostra, não de forma linear ao longo de toda a distribuição etária (Seção~\ref{sec:r_h3}). E o tamanho do pacote negociado acompanha, de fato, a severidade da lacuna que a equipe adquirente buscava resolver. É a correlação mais forte entre as quatro hipóteses centrais, ainda que classificada como efeito fraco (Seção~\ref{sec:r_h4}). Vale registrar também o resultado complementar sobre o objetivo sazonal da equipe: a leitura inicial, de que times em reconstrução (\textit{tanking}) melhoram de forma inesperada depois de uma troca, não resiste a um exame mais cuidadoso. A mesma melhora aparece, em grau semelhante, nos três perfis de objetivo sazonal testados, e a correlação negativa entre o desempenho pré-troca e o ganho obtido depois dela, presente nos três grupos, é a assinatura clássica de regressão à média, não de um efeito específico de estratégia competitiva (Seção~\ref{sec:r_context}). Ignorar essa checagem teria produzido uma manchete mais interessante e uma conclusão errada.

Em todos os seus quatro testes centrais, este é um estudo correlacional sobre dados observacionais, não um experimento. Por isso, nenhum resultado relatado no Capítulo 4 foi aceito como está sem antes passar por uma verificação de robustez, uma análise por subgrupo ou uma checagem de explicação alternativa.

\section{Limitações do estudo}
\label{sec:final_limitations}

A principal limitação deste trabalho é de natureza estrutural, uma vez que nenhum dos resultados aqui relatados identifica mecanismo causal. A correlação entre a variação da desigualdade salarial e a variação de desempenho, por exemplo, é compatível com a Teoria da Equidade, mas também é compatível com uma explicação reversa: equipes que já vinham jogando melhor podem ser, por outros motivos, as mesmas que optam por reter seus jogadores mais bem pagos em vez de trocá-los. Também é compatível com uma terceira variável não observada que influencia as duas ao mesmo tempo, como lesões, mudanças de comissão técnica ou a força específica do calendário remanescente, esta última apenas parcialmente controlada pela Performance Ajustada. O mesmo vale, com a devida adaptação, para as outras três hipóteses. Delinear esse tipo de correlação com rigor estatístico, como foi feito aqui, é uma etapa necessária antes de qualquer afirmação causal, mas não é a mesma coisa que fazê-la.

Quatro limitações de natureza mais operacional, já registradas no Capítulo 3 (Seção~\ref{sec:method_limitations}) no momento em que a metodologia foi definida, se confirmaram como relevantes ao longo da análise dos resultados. A heterogeneidade da fonte do WPA entre 2012-13 e 2024-25 deixou de ser uma preocupação abstrata e se tornou, na prática, a explicação mais plausível para a instabilidade do resultado contínuo entre eras. A cobertura parcial do arquétipo competitivo pré-troca, hoje disponível apenas para as safras de 2018-19 a 2020-21, restringiu a leitura descritiva complementar de trocas por perfil de equipe a um recorte pequeno da amostra total de eventos. E o poder estatístico reduzido de dividir a amostra do objetivo sazonal em três grupos, antecipado desde a Seção~\ref{sec:method_limitations}, exigiu cautela redobrada ao interpretar tanto o resultado nulo dos \textit{contenders} quanto o resultado positivo, e inesperado, do \textit{tanking}. Só a checagem de regressão à média (Seção~\ref{sec:r_context}) confirmou que essa cautela era, de fato, necessária. Por fim, como o arquétipo competitivo usado para moderar a primeira hipótese (Subseção~\ref{sec:r_h1_moderation}) é calculado sobre a temporada regular completa de cada time, sem isolar os jogos anteriores de uma eventual troca dos posteriores a ela, existe uma possibilidade teórica de vazamento temporal entre essa variável moderadora e o próprio desempenho pós-troca que ela ajuda a explicar, um efeito que este trabalho não isola.

Há também uma limitação de consistência interna que vale reconhecer explicitamente: a verificação de que o efeito da desigualdade salarial sobre o desempenho se concentra no arquétipo Ofensivo (Subseção~\ref{sec:r_h1_moderation}) usa o \textit{Simple Rating System} como medida de desempenho, e não a Performance Ajustada empregada no teste central da mesma hipótese, porque o recálculo do arquétipo de cada time-temporada depende de uma consulta direta ao banco de dados que não carrega consigo o ajuste pela força do adversário. O achado sobre o arquétipo Ofensivo é, por isso, indicativo, e soma-se a um segundo problema, este estatístico: com quatro arquétipos testados independentemente a $\alpha = 0{,}05$, a chance de pelo menos um resultado "significativo" aparecer por acaso já é de cerca de 18,5\%, e nenhuma correção para múltiplas comparações foi aplicada nesse teste exploratório específico.

A terceira hipótese carrega uma limitação de validade de construto que também merece registro explícito: a idade do jogador adquirido foi usada como \textit{proxy} para sua experiência profissional na NBA (Seção~\ref{sec:r_h3}), mas as duas coisas não são necessariamente equivalentes. Um jogador de 29 anos que estreou tardiamente na liga pode somar menos temporadas de NBA, e por extensão menos experiência efetiva no nível de competição testado, do que outro jogador de 26 anos que estreou aos 19 e já acumula oito temporadas profissionais. Como a base de dados usada neste trabalho registra a idade do jogador no momento da troca, mas não o número de temporadas de NBA efetivamente jogadas por ele, essa distinção não pôde ser testada diretamente, e o resultado da Hipótese 3 deve ser lido como um efeito de idade cronológica, não estritamente de experiência de carreira.

Por fim, a amostra usada nos testes centrais é toda extraída de uma única liga, a NBA, sob uma única definição de janela pré-\textit{trade deadline} (30 dias, testada como robusta entre 14 e 60 dias na Seção~\ref{sec:r_h1_main}) e cobrindo treze temporadas específicas, marcadas por mudanças de regras de teto salarial e de disponibilidade de dado que não foram tratadas como variáveis do modelo. Se as associações aqui relatadas se sustentam em outras ligas de basquete, sob outras estruturas de teto salarial, ou mesmo em temporadas futuras da própria NBA sob um novo acordo coletivo, é uma pergunta que este trabalho não tem como responder.

\section{Implicações gerenciais}
\label{sec:managerial_implications}

Para um \textit{front office}, a implicação mais direta sobre a desigualdade salarial não é "evite concentrar a folha salarial", em termos gerais. O efeito é real, mas pequeno demais para justificar, sozinho, recusar uma troca vantajosa sob outros critérios. A implicação mais útil é mais específica: o resultado sugere que o custo de aumentar a concentração salarial em quadra não é uniforme entre todos os elencos, e parece pesar mais especificamente sobre equipes de perfil ofensivo e ritmo acelerado. Uma equipe desse perfil que esteja avaliando concentrar sua folha salarial em uma ou duas estrelas, às custas do restante da rotação, tem, segundo este resultado, um motivo a mais para essa cautela do que uma equipe de perfil defensivo, na qual a mesma reorganização salarial não mostrou associação estatística com o desempenho.

O ajuste funcional do jogador adquirido deixa, na prática, uma lição sobre os limites de uma métrica, não sobre sua própria irrelevância. A lógica de que um jogador que preenche uma deficiência real do elenco tende a valer mais do que um jogador redundante continua defensável em termos de basquete, e a literatura do \textit{O-ring} de \citet{kremer1993} segue oferecendo um argumento teórico sólido nesse sentido. O que este trabalho mostra é que operacionalizar essa lógica de forma automática e categórica, a partir de um corte único da Contribuição Relativa, ainda não produz um sinal confiável o suficiente para orientar decisões sozinho. Um analista de \textit{front office} que use essa métrica deveria tratá-la como mais um insumo qualitativo entre vários, e não como um critério decisivo de avaliação de alvo de troca, ao menos até que a instabilidade entre eras de dado, discutida na Seção~\ref{sec:final_limitations}, seja melhor compreendida.

Já a experiência dos jogadores adquiridos e a estrutura da negociação sugerem heurísticas mais diretamente aplicáveis, com a devida cautela sobre o tamanho modesto dos efeitos. Times buscando um reforço com impacto imediato, tipicamente equipes competindo por vaga de \textit{playoff} no curto prazo, encontram nos dados deste trabalho um argumento a favor de priorizar experiência sobre potencial, especialmente ao considerar jogadores no quartil mais velho da distribuição etária de mercado. E times avaliando o custo de uma troca deveriam esperar, com alguma confiança, que preencher uma lacuna genuinamente severa do elenco tende a exigir movimentar mais de um jogador. Isso serve tanto como orientação de planejamento quanto como sinal de alerta: uma proposta de troca que promete resolver uma carência importante do time com um único jogador de baixo custo deveria, pelos padrões observados nesta amostra, ser recebida com ceticismo adicional.

Por fim, o achado sobre o objetivo sazonal (Seção~\ref{sec:r_context}) carrega uma implicação menos técnica, mas igualmente prática: gestores e comentaristas costumam interpretar a melhora de desempenho de um time em reconstrução logo após uma troca como evidência de que a estratégia de \textit{tanking} "funcionou" competitivamente no curto prazo. Este trabalho encontra o mesmo padrão de melhora nos três perfis de objetivo sazonal testados, associado ao ponto de partida de cada equipe, não ao objetivo declarado dela. Isso é, em si, um alerta útil contra atribuir a uma decisão estratégica um efeito que a estatística já explica por outro caminho.

\section{Sugestões para trabalhos futuros}
\label{sec:future_work}

A extensão mais imediata, e a mais viável do ponto de vista de engenharia de dados, é completar o cruzamento entre a base de trocas e a clusterização de arquétipos competitivos para as dez temporadas hoje não cobertas por essa marcação (Seção~\ref{sec:final_limitations}), o que permitiria repetir a análise por perfil de equipe da Subseção~\ref{sec:r_h1_moderation} usando a própria Performance Ajustada, em vez do \textit{Simple Rating System}, e assim testar se o efeito concentrado no arquétipo Ofensivo se sustenta com a variável de desempenho correta e sob correção para múltiplas comparações.

Uma segunda linha, mais custosa, mas potencialmente mais reveladora, é buscar uma fonte de dado posicional jogo a jogo que cubra as treze temporadas de forma homogênea, eliminando a fronteira de qualidade entre as três eras de disponibilidade de dado já descritas na Seção~\ref{sec:h2_wpa_ops} do Capítulo 3 (Tabela~\ref{tab:h2_eras}), identificada como a explicação mais provável para a instabilidade do resultado contínuo. Sem essa homogeneização, não é possível saber se o ajuste funcional do jogador adquirido é, de fato, um preditor fraco de desempenho, ou se apenas parece fraco por causa do dado disponível.

Do ponto de vista teórico, duas extensões parecem particularmente promissoras. A primeira é testar se a relação entre desigualdade salarial e desempenho encontrada neste trabalho é de fato linear, ou se segue o padrão em U invertido que \citet{ching2019} relatou usando dados agregados por temporada. Este estudo não testou formalmente um termo quadrático, e a amostra de eventos de troca, mais granular do que a usada por \citet{ching2019}, poderia revelar um ponto de inflexão que a análise linear aqui empregada não captura. A segunda é aprofundar por que esse efeito se concentra especificamente no arquétipo Ofensivo. É uma pergunta que dados quantitativos de composição de elenco, como os usados aqui, não respondem sozinhos, e que poderia se beneficiar de fontes complementares, como entrevistas com membros de comissão técnica ou medidas de coesão de vestiário.

Por fim, como toda a base de dados deste trabalho vem de uma única liga, replicar o mesmo modelo multidimensional em outras ligas de basquete profissional, como a WNBA ou ligas europeias, cada uma sob sua própria estrutura de teto salarial e seu próprio mercado de trocas, permitiria testar se as associações aqui relatadas são uma particularidade do arranjo institucional da NBA ou um padrão mais geral do esporte coletivo profissional. Enquanto isso não acontece, os resultados deste trabalho devem ser lidos como o que são: evidência específica de uma liga, de um recorte temporal e de um conjunto de métricas, consistente o suficiente para ser levada a sério, mas modesta demais para ser tratada como resposta definitiva.
